% Generated draft: requires mathematical and source review.
\documentclass[11pt]{article}
\usepackage[margin=25mm]{geometry}
\usepackage{fontspec}
\setmainfont{Latin Modern Roman}
\usepackage{amsmath,amssymb,mathtools,unicode-math}
\setmathfont{Latin Modern Math}
\usepackage{enumitem,needspace,xurl,hyperref}
\setlength{\parindent}{0pt}
\setlength{\parskip}{5pt}
\setlength{\emergencystretch}{3em}
\newcommand{\N}{\mathbb{N}}
\newcommand{\Z}{\mathbb{Z}}
\newcommand{\Q}{\mathbb{Q}}
\newcommand{\R}{\mathbb{R}}
\newcommand{\C}{\mathbb{C}}
\newcommand{\F}{\mathbb{F}}
\newcommand{\A}{\mathbb{A}}
\newcommand{\PP}{\mathbb{P}}
\newcommand{\E}{\mathbb{E}}
\newcommand{\eps}{\varepsilon}
\newcommand{\dd}{\,\mathrm{d}}
\begin{document}
\section*{UnsolvedMath ID: 800395}
\textbf{Source problem code:} catalogue entry
\subsection{Definitions and mathematical statement}
Let \(T\) be a nonsingular algebraic curve, and consider the family written in the supplied source as
\[
T\longrightarrow\mathcal M_d^{\mathrm{crit}},\qquad t\longmapsto(f_t,P_{1,t},\ldots,P_{2d-2,t}).
\]
Here \(\mathcal M_d^{\mathrm{crit}}\) denotes the source's moduli space of degree-\(d\) rational maps with marked critical points, and \(P_{i,t}\) is the critical point bearing label \(i\) in the fiber over \(t\). The source does not specify the ground field, the allowed range of \(d\), or the convention for repeated critical points. These qualifications remain unresolved.

A rational map \(g\) is \emph{postcritically finite} (PCF) if every critical point \(c\) has finite forward orbit:
\[
\#\{g^n(c):n\in\Z_{\geq0}\}<\infty.
\]
The supplied conjecture, attributed to Baker--DeMarco (2013), asserts the equivalence of the following conditions:
\begin{enumerate}
\item[(a)] The set \(\{t\in T:f_t\text{ is PCF}\}\) is infinite.
\item[(b)] For every \(i,j\in\{1,\ldots,2d-2\}\), the marked critical-point maps \(P_i:T\to\PP^1\) and \(P_j:T\to\PP^1\) are \(f\)-dynamically related.
\end{enumerate}
The quantification in (b) includes \(i=j\), as written in the source. The supplied record does not define ``\(f\)-dynamically related,'' so condition (b) is retained as source terminology and does not yet give a self-contained mathematical condition. The field over which parameters are counted in (a) is also unspecified. Furthermore, the displayed morphism to a moduli space does not by itself specify a global representative family on a fixed \(\PP^1\); the source nevertheless writes the critical markings as maps to \(\PP^1\). The intended representative-family or lifting convention requires clarification.
\subsection{Short English statement}
The supplied conjecture says that a one-dimensional family of rational maps with marked critical points has infinitely many postcritically finite members exactly when every pair of marked critical points is dynamically related. Precision requires the missing definition of dynamical relatedness, the ground field and parameter domain, and the conventions for the marked family.
\subsection{Sources}
\begin{itemize}
\item Local UnsolvedMath registry: permanent ID \(800395\), source problem code \texttt{catalogue entry}, in the supplied AMR Open Problem Lists record. The registry labels the problem \emph{open}; this is supplied classification, not an independently verified status assessment.
\item Benedetto, DeMarco, Ingram, Jones et al., \emph{Current Trends and Open Problems in Arithmetic Dynamics} (2018), Conjecture 18, identified in the supplied record as the Dynamical Andr\'{e}--Oort Conjecture for 1-Dimensional Families. Supplied source URL: \url{https://arxiv.org/abs/1806.04980}. The record reports access on 2026-10-16 and extraction from arXiv LaTeX; these provenance details are supplied and have not been independently verified here.
\item The extracted statement attributes the conjecture to Baker--DeMarco (2013), with reference label \texttt{[arxiv1211.0255]}. No further bibliographic details for that reference are supplied.
\end{itemize}
\end{document}
